\documentclass[]{article}
\usepackage{lmodern}
\usepackage{amssymb,amsmath}
\usepackage{ifxetex,ifluatex}
\usepackage{fixltx2e} % provides \textsubscript
\ifnum 0\ifxetex 1\fi\ifluatex 1\fi=0 % if pdftex
  \usepackage[T1]{fontenc}
  \usepackage[utf8]{inputenc}
\else % if luatex or xelatex
  \ifxetex
    \usepackage{mathspec}
  \else
    \usepackage{fontspec}
  \fi
  \defaultfontfeatures{Ligatures=TeX,Scale=MatchLowercase}
\fi
% use upquote if available, for straight quotes in verbatim environments
\IfFileExists{upquote.sty}{\usepackage{upquote}}{}
% use microtype if available
\IfFileExists{microtype.sty}{%
\usepackage[]{microtype}
\UseMicrotypeSet[protrusion]{basicmath} % disable protrusion for tt fonts
}{}
\PassOptionsToPackage{hyphens}{url} % url is loaded by hyperref
\usepackage[unicode=true]{hyperref}
\hypersetup{
            pdftitle={Language, Cognition, and the Brain},
            pdfauthor={Ethan Weed},
            pdfborder={0 0 0},
            breaklinks=true}
\urlstyle{same}  % don't use monospace font for urls
\usepackage[margin=1in]{geometry}
\usepackage{natbib}
\bibliographystyle{plainnat}
\usepackage{longtable,booktabs}
% Fix footnotes in tables (requires footnote package)
\IfFileExists{footnote.sty}{\usepackage{footnote}\makesavenoteenv{long table}}{}
\usepackage{graphicx,grffile}
\makeatletter
\def\maxwidth{\ifdim\Gin@nat@width>\linewidth\linewidth\else\Gin@nat@width\fi}
\def\maxheight{\ifdim\Gin@nat@height>\textheight\textheight\else\Gin@nat@height\fi}
\makeatother
% Scale images if necessary, so that they will not overflow the page
% margins by default, and it is still possible to overwrite the defaults
% using explicit options in \includegraphics[width, height, ...]{}
\setkeys{Gin}{width=\maxwidth,height=\maxheight,keepaspectratio}
\IfFileExists{parskip.sty}{%
\usepackage{parskip}
}{% else
\setlength{\parindent}{0pt}
\setlength{\parskip}{6pt plus 2pt minus 1pt}
}
\setlength{\emergencystretch}{3em}  % prevent overfull lines
\providecommand{\tightlist}{%
  \setlength{\itemsep}{0pt}\setlength{\parskip}{0pt}}
\setcounter{secnumdepth}{5}
% Redefines (sub)paragraphs to behave more like sections
\ifx\paragraph\undefined\else
\let\oldparagraph\paragraph
\renewcommand{\paragraph}[1]{\oldparagraph{#1}\mbox{}}
\fi
\ifx\subparagraph\undefined\else
\let\oldsubparagraph\subparagraph
\renewcommand{\subparagraph}[1]{\oldsubparagraph{#1}\mbox{}}
\fi

% set default figure placement to htbp
\makeatletter
\def\fps@figure{htbp}
\makeatother

\usepackage{booktabs}
\usepackage{amsthm}
\makeatletter
\def\thm@space@setup{%
  \thm@preskip=8pt plus 2pt minus 4pt
  \thm@postskip=\thm@preskip
}
\makeatother
\usepackage{etoolbox}
\makeatletter
\providecommand{\subtitle}[1]{% add subtitle to \maketitle
  \apptocmd{\@title}{\par {\large #1 \par}}{}{}
}
\makeatother

\title{Language, Cognition, and the Brain}
\providecommand{\subtitle}[1]{}
\subtitle{An Introductory Course in Psycholinguistics at Aarhus University}
\author{Ethan Weed}
\date{17 April, 2020}

\begin{document}
\maketitle

{
\setcounter{tocdepth}{2}
\tableofcontents
}
\section{Introduction}\label{introduction}

\subsection{What is Language?}\label{what-is-language}

What is language? This is one of those questions that seems like it has
an easy and obvious answer, but like most questions, the more I think
about it, the harder it becomes to give a single, satisfactory answer.
Language is something we use every day, so of course we know what it is.
It's how we communicate with each other. Then again, it is not the
\emph{only} way we communicate with each other. If I go to the
supermarket, and another shopper and I move to give space to each other
in the aisle, we are intentionally communicating information about how
we can move our bodies and shopping carts to allow us both to get the
items we want from the shelves. But this doesn't really seem like an
example of language. If the same fellow shopper and I misjudge each
other's intentions, and move the wrong way, and smile to each other,
this is also communication, but it still doesn't feel like language. You
can easily think of other examples of humans communicating with other
humans in ways that we wouldn't normally call language: music, painting,
photography, dance\ldots{}

\section{Classic Models of Aphasia}\label{classic-models-of-aphasia}

\section{Primary progressive aphasia}\label{primary-progressive-aphasia}

\section{Speech perception and the dual-stream
model}\label{speech-perception-and-the-dual-stream-model}

\section{More on Speech Perception and the Dual Stream
Model}\label{more-on-speech-perception-and-the-dual-stream-model}

\section{Speech Production}\label{speech-production}

\section{Prosody}\label{prosody}

\section{Word Meaning}\label{word-meaning}

\section{Action Verbs}\label{action-verbs}

\section{Reading and Spelling}\label{reading-and-spelling}

\section{Developmental Dyslexia}\label{developmental-dyslexia}

\section{Right Hemisphere Disorder}\label{right-hemisphere-disorder}

\section{Introduction}\label{introduction-1}

\subsection{What is Language?}\label{what-is-language-1}

What is language? This is one of those questions that seems like it has
an easy and obvious answer, but like most questions, the more I think
about it, the harder it becomes to give a single, satisfactory answer.
Language is something we use every day, so of course we know what it is.
It's how we communicate with each other. Then again, it is not the
\textbf{only} way we communicate with each other. If I go to the
supermarket, and another shopper and I move to give space to each other
in the aisle, we are intentionally communicating information about how
we can move our bodies and shopping carts to allow us both to get the
items we want from the shelves. But this doesn't really seem like an
example of language. If the same fellow shopper and I misjudge each
other's intentions, and move the wrong way, and smile to each other,
this is also communication, but it still doesn't feel like language. You
can easily think of other examples of humans communicating with other
humans in ways that we wouldn't normally call language: music, painting,
photography, dance\ldots{}

\section{Classic Models of Aphasia}\label{classic-models-of-aphasia-1}

\section{Primary progressive
aphasia}\label{primary-progressive-aphasia-1}

\section{Speech perception and the dual-stream
model}\label{speech-perception-and-the-dual-stream-model-1}

\section{More on Speech Perception and the Dual Stream
Model}\label{more-on-speech-perception-and-the-dual-stream-model-1}

\section{Speech Production}\label{speech-production-1}

\section{Prosody}\label{prosody-1}

\section{Word Meaning}\label{word-meaning-1}

\section{Action Verbs}\label{action-verbs-1}

\section{Reading and Spelling}\label{reading-and-spelling-1}

\section{Developmental Dyslexia}\label{developmental-dyslexia-1}

\section{Right Hemisphere Disorder}\label{right-hemisphere-disorder-1}

\end{document}
